\section{数据采集与 Replay Buffer 管理}

\subsection{好数据塑造通用性}

Q-Learning 是典型的 off-policy 方法，训练过程中可重用历史转移，因此数据质量直接决定效果。

\begin{importantbox}{数据的三维质量框架}
\begin{enumerate}
    \item \textbf{Coverage}：状态-动作对是否遍布整个训练分布。
    \item \textbf{Freshness}：旧数据可能已经过时，尤其当策略迭代很快时。
    \item \textbf{Signal-to-Noise}：奖励与下一状态之间的因果链是否足够清晰。
\end{enumerate}
\end{importantbox}

\begin{figure}[H]
\centering
\includegraphics[width=0.86\textwidth,page=1]{06_cs224r_qlearning_2025.pdf}
\caption{Slide：Replay buffer 抽样的分布偏移与纠正策略}
\end{figure}

\subsection{优先经验回放与分布修正}

优先经验回放（PER）根据 TD Error 分配采样概率，使训练更多聚焦于高价值/高误差样本。修正方法包括：

\begin{itemize}
    \item importance sampling weight 用以抵消非均匀采样
    \item alpha、beta 超参数控制采样与修正的强度
\end{itemize}

\begin{knowledgebox}{Replay Buffer 的工程实践提示}
1. 定期清理过旧元素以防止“冷启动”偏移。 2. 保持 buffer size 在百万级别以平衡样本多样性。 3. 为不同任务维护多个 buffer（exploration vs. demonstration）。
\end{knowledgebox}

\subsection{案例：Atari Replay 抽样流程}

以经典 Atari Q-Learning 为例，训练流水线可简化为以下步骤：

\begin{enumerate}
    \item 通过 epsilon-greedy 收集原始 transition，插入 `online buffer`。
    \item 每 1k steps 用 prioritized sampling 生成 mini-batch，同时记录 TD error distribution。
    \item 将 TD error 作曲线可视化，若长尾方向持续偏移，触发 `replay augmentation`（#32, #64 sampling）。
    \item 定期将 high-variance batch 提前放入 `high-priority buffer` 供评估与 debugging。
\end{enumerate}

\begin{importantbox}{Replay 监控 checklist}
\begin{itemize}
    \item 记录 replay hit ratio，确保采样重用率在 70-95%。
    \item `TD error` 上升时自动增加 `beta` 以增强 importance sampling 校正。
    \item 将 human-verified failure transitions 复制到 separate buffer，方便 later inspection。
\end{itemize}
\end{importantbox}

\subsection{本章小结}

Q-Learning 的 off-policy 特性让 replay buffer 成为性能发动机，采样策略和数据质量直接决定最终泛化能力。

